\documentclass[aspectratio=169]{beamer}


\usepackage{draftwatermark}
\setbeamercolor{background canvas}{bg=}%transparent canvas


\usepackage{graphicx}

\definecolor{egyptianblue}{rgb}{0.06, 0.2, 0.65}
\usepackage{xcolor}

\setbeamercovered{transparent}
%\usepackage{pxfonts}
\usepackage{listings}
\lstset{language=bash,
                basicstyle=\ttfamily,
                keywordstyle=\color{cyan}\ttfamily,
                stringstyle=\color{orange}\ttfamily,
                commentstyle=\color{gray}\ttfamily,
                morecomment=[l][\color{magenta}]{\#}
}

\lstdefinestyle{BashInputStyle}{
  language=bash,
  basicstyle=\small\sffamily,
  numbers=left,
  numberstyle=\tiny,
  numbersep=3pt,
  frame=tb,
  columns=fullflexible,
  backgroundcolor=\color{yellow!20},
  linewidth=0.9\linewidth,
  xleftmargin=0.1\linewidth
}
\sffamily

\usepackage{tipa}

\usepackage{natbib}

\titlegraphic{\includegraphics[height=.2\textheight]{../hidrus-logo.pdf}}

\begin{document}
%--------------------------------------------------------------------------
%title page
%--------------------------------------------------------------------------
\title{\Large \noindent \textbf{\textcolor{egyptianblue}{Notes on \underline{Hi}gh-resolution \underline{D}ownscaling of \underline{R}ainfall \underline{U}sing \underline{S}TEPS (HiDRUS)}}}  

\author{\emph{\href{www.baraut.info}{Bhupendra A Raut}}} 
\date{\today} 

\frame{\titlepage} 
\frame{\frametitle{Table of contents}\tableofcontents} 

%--------------------------------------------------------------------------
%Introduction
%--------------------------------------------------------------------------
\section{Introduction}
%
\begin{frame}{STEPS} %STEPS
The  \underline{S}hort \underline{T}erm \underline{E}nsemble \underline{P}rediction \underline{S}ystem (STEPS)  has several modules to 

\begin{itemize}
\item disaggregate rainfall \citep{seed2013steps}
\item simulate time-series of mean area rainfall \citep{seed2000BL}
\item make design-storm simulations \citep{seed1999steps} 
\item and to produce seamless forecast when cascaded with NWP models \citep{bowler2006steps}
\end{itemize}

STEPS is mainly used for nowcasting and multiple simulations of the past extreme events.
\end{frame}
%
%
%
%
%
\begin{frame}{HiDRUS} %HiDRUS
\begin{itemize}
\item HiDRUS\footnote{Pronounced ``\textipa{ha\;Idr@s}''} is an implementation of the STEPS library to downscale GCM/RCM rainfall to a very high space-time resolution (1 km, 6 minutes).
\item HiDRUS-1 is a fully stochastic downscaling scheme which uses a broken-line model to generate the mean area rainfall time-series (MARTS). Linear regression models are use to generate cascade parameters for each time-step. 
\item HiDRUS-2 uses sampling of the past MARTS to generate future rainfall projections. 
\end{itemize}
This document describes the steps in running HiDRUS-2 model to downscale rainfall.
\end{frame}
%
%%
%
%
%
\section{Running HiDRUS-2}
%
\begin{frame}{Running HiDRUS-2}
\begin{columns}
\column{0.6\textwidth}
HiDRUS-2 methodology has two configurations.
\begin{enumerate}
\item \textbf{perfect mode (A)} to downscale radar data with computed parameters. This is used for model validation, model bias \& uncertainty estimation.
\item \textbf{projection mode (B)} to be run with GCM/RCM/reanalysis data.
\end{enumerate}
The flow-digram shows important stages of the \emph{perfect} and the \emph{projection} modes of HiDRUS-2 methodology.

\column{0.5\textwidth}
	\includegraphics[height=3in]{../fig/flow.pdf}
\end{columns}
\end{frame}
%
%
%
%
%
%--------------------------------------------------------------------------
%Compute parameters
%--------------------------------------------------------------------------
\begin{frame}[fragile]
\frametitle{Compute Cascade Parameters}
\begin{columns}
\column{0.6\textwidth}
\begin{lstlisting}[style=BashInputStyle]
$mkcascparm <fPath> <rn>
\end{lstlisting}

fPath - path of the input directory,\\
rn - Radar id. This will be suffixed to the output file name.

\begin{lstlisting}[style=BashInputStyle]
$./mkcascparm ~/data/ MLB
$cd ./c_parms_dbz
$ls
 20100108_MLB.prm
 ...
 ...
 20101224_MLB.prm
 ...
 RainPMatch_MLB.ima  
 RainPMatch_MLB.ind  
 RainPMatch_MLB.pjn
 ...
\end{lstlisting}

\column{0.4\textwidth}
	\includegraphics[height=2in]{../fig/casc_prm.pdf}
\end{columns}

\end{frame}
%
%
%
%--------------------------------------------------------------------------
%Perfect Run
%--------------------------------------------------------------------------
%
\section{Perfect Run}

\begin{frame}[fragile]
\frametitle{Perfect Run - Make Parameters}

\begin{columns}

\column{0.4\textwidth}
\includegraphics[height=2.5in, width=2in]{../fig/perf_prm.pdf}

\column{0.6\textwidth}
\begin{lstlisting}[style=BashInputStyle] 
$vim mk-HiDRUS2-Parms_pppr.R

#set parameters directory
setwd("~/data/c_parms_dbz/")

#input file name pattern
fn_pat <- "_MLB.prm"  

#desired output file name
ofName<-"simParm-2010.txt"
\end{lstlisting}

Run the script.
\begin{lstlisting}[style=BashInputStyle]
$Rscript mk-HiDRUS2-Parms_pppr.R 
\end{lstlisting}
\end{columns}

\end{frame}
%
%
%
%
%
%
%
\begin{frame}[fragile]
\frametitle{Perfect Run - Config and Run}

\begin{columns}
\column{0.4\textwidth}
	\includegraphics[width=2.5in]{../fig/runPerf_h2.pdf}

\column{0.6\textwidth}
Provide appropriate paths in config.txt file
\begin{lstlisting}[style=BashInputStyle]
$vim config.txt

outFilePathPrefix = /path/pSim2010 ;
pMatchPathPrefix = /path/RainPMatch_MLB ;
prmFileName = /path/simParm-2010.txt ;
\end{lstlisting}


\begin{lstlisting}[style=BashInputStyle] 
$hidrus2 ./config.txt ens_id
\end{lstlisting}
If you give ens\_id here then ens\_id in config.txt will be overridden. This feature can be used to make several simulations with the same config file.

\end{columns}
\end{frame}
%
%
%
%
%
%
\begin{frame}[fragile]
\frametitle{Perfect Run - Excerpt of a PBS Job Script}
\begin{lstlisting}
#!/bin/bash
# Example submission of PBS jobs in a loop
NUMBERS=$(seq 1 100) 
for NUM in ${NUMBERS}
do
  NAME=h2pp${NUM}
  echo "Submitting: ${NAME}"	
  PBS="#!/bin/bash\n\
  #PBS -N ${NAME}\n\
  #PBS -l walltime=10:00:00\n\
  #PBS -l mem=500MB\n\
  #PBS -l wd\n\
  ./hidrus2 ./config.txt ${NUM}"
	echo -e ${PBS} | qsub 
	echo "done."
done

\end{lstlisting}

\end{frame}
%
%
%
%
%
%
\begin{frame}[fragile]
\frametitle{Perfect Run - Output}
In output directory
\begin{lstlisting}[style=BashInputStyle]
$ls
pSim2010_001_tStat.asc  pSim2010_055_tStat.asc
pSim2010_002_tStat.asc  pSim2010_056_tStat.asc
...
pSim2010_001.nc  pSim2010_055.nc
pSim2010_002.nc  pSim2010_056.nc
...
\end{lstlisting}

The files ``$\ast\_tStat.asc$'' contains domain average properties for all rainy time-steps.
\begin{lstlisting}[style=BashInputStyle]
$head pSim2010_055_tStat.asc
time	 pm	in_dbz_fmean	in_dbz_fstd	in_rain_fmean	in_rain_fstd	in_rain_cmean
in_rFrac	in_east	in_south	out_dbz_fmean	out_dbz_fstd	out_rain_fmean	
out_rain_fstd	out_rain_cmean	out_rFrac
1200018240	1200018273	1.83	 7.65	0.52	 3.93	8.6 3	0.06	 2.0	 0.2 	
1.83	 7.64	0.52 	3.92	 8.62 	0.06
\end{lstlisting}

\end{frame}
%
%
%
%
%
\begin{frame}[fragile]
\frametitle{Perfect Run - Output}
\begin{lstlisting}[style=BashInputStyle]
$ncdump -h pSim2010_055.nc 
netcdf pSim2010_055 {
dimensions:
	lat = 230 ;
	lon = 256 ;
	time = UNLIMITED ; // (23000 currently)
variables:
	float lat(lat) ;
    lat:units = "degrees_north" ;
	float lon(lon) ;
    lon:units = "degrees_east" ;
	int64 time(time) ;
    time:units = "seconds since 1970-01-01 00:00:00 UTC" ;
	float rain(time, lat, lon) ;
    rain:units = "mm/hr" ;
    rain:_FillValue = -999.f ;
    rain:long_name = "instantaneous rainrate" ;
    rain:_description = "perfect_simulations" ;
// global attributes:
		:HiDRUS = "developed at Monash University and Bureau of Meteorology." ;
		:creator_name = "Bhupendra A. Raut and Alan W. Seed" ;
		:creator_webpage = "www.baraut.info" ;
		:_institution = "School of Earth Atmosphere and Environment, Monash University" ;
		:_acknowledgment = "Bureau of Meteorology and CRC for Water Sensitive Cities" ;
		:_project = "B1.1: Cities as Water Supply Catchments – Urban Rainfall in a changing climate" ;
		:_project_leaders = "Christian Jakob, Michael J. Reeder" ;
		:date_created = "Tue Sep 29 11:53:59 2015" ;
}

\end{lstlisting}


\end{frame}
%
%
%
%
%
\section{Projection Run}
\begin{frame}[fragile]
\frametitle{Projection Run - Make Rain Events Library}
\begin{lstlisting}[style=BashInputStyle]
$vim mk_rainEvents_lib.R
\end{lstlisting}

\begin{lstlisting}[language=R]
#set working directory to parameters directory
setwd("/home/bhupendra/data/c_parms_dbz/")
ofileName<-"RainLib_2010_MLB"  #no extension

#--------- Some initialisation settings----------------#
window <- 240   #number of point in sampled ts
cut_off=0.1     #mean area rain mm/hr
duration=5      #min duration of rain in tSteps
nprms=65        # number of parameters 
fn_pat <- "_MLB.prm"    #file name pattern 
\end{lstlisting}

\end{frame}



\begin{frame}[fragile]
\frametitle{Projection Run - Make Rain Events Library}
\begin{columns}

\column{0.6\textwidth}
\begin{lstlisting}[style=BashInputStyle] 
$Rscript mk_rainEvents_lib.R
$ls 
RainLib_2010_MLB.nc RainLib_2010_MLB.pdf
\end{lstlisting}
The cascade parameters for each event will be stored in the NetCDF file.
\begin{lstlisting}[style=BashInputStyle] 
$ncdump -h RainLib_2010_MLB.nc
dimensions:
    id = UNLIMITED ; // (1546 currently)
    prms = 65 ;
    tstep = 240 ;
variables:
    int id(id) ;
    id:long_name = "identity number 
    for the event" ;
	double ts_prms(id, tstep, prms) ;
    ts_prms:long_name = "cascade parameters" ;
\end{lstlisting}

\column{0.5\textwidth}
\includegraphics[width=2.8in]{../fig/mkRainlib.pdf}

\end{columns}
\end{frame}


\begin{frame}[fragile]
\frametitle{Projection Run - Make Synoptic Conditions Library}
\begin{columns}

\column{0.6\textwidth}
At this stage auxiliary data can be stored for the back-identification of events in future projections.
This may include 
\begin{itemize}
\item Weather Regimes
\item vertical profiles of U, V, humidity etc.
\item MSLP patterns, 
\item accumulation of rain from Rain-gauge, TRMM
\end{itemize}

\column{0.5\textwidth}
\includegraphics[width=2.8in]{../fig/mkRainlib.pdf}

\end{columns}
\end{frame}

\begin{frame}[fragile]
\frametitle{Projection Run - Select Future Parameters}
\begin{columns}

\column{0.45\textwidth}
\begin{itemize}
\item An example script mk-HiDRUS2-Parms.R is provided for selecting future days.
\item Users will have to use their matching criterion for the past and future days.
\end{itemize}

\begin{lstlisting}[style=BashInputStyle] 
$cdo fldmean -sellonlatbox,143.5,146,-36.5,-39 ifile.nc ofile.nc
\end{lstlisting}


\column{0.6\textwidth}
\includegraphics[width=3.5in]{../fig/selFuturePrm.pdf}

\end{columns}
\end{frame}

\begin{frame}[fragile]
\frametitle{Projection Run - Config and Run}
\begin{columns}

\column{0.52\textwidth}
\begin{itemize}
\item Running HiDRUS-2 for projection is same as for the perfect run.
\item An optional argument ``parmFilePath''  can be used to provide parameters file path. This will over-write the path of the parameter file in config.txt
\item This allows us to run HiDRUS-2 for several input parameter files without modifying the config file.
\end{itemize}
possible run commands
\begin{lstlisting}[style=BashInputStyle]
$hidrus2 ./config.txt 
$hidrus2 ./config.txt ens_id
$hidrus2 ./config.txt ens_id parmFilePath
\end{lstlisting}

\column{0.5\textwidth}
\includegraphics[width=2.9in]{../fig/runH2.pdf}

\end{columns}
\end{frame}
%
%
%
%
\begin{frame}[fragile]
\frametitle{Projection Run - PBS Job Script}
\begin{lstlisting}[language=Python]
from popen2 import popen2
import time

yesno='no'
yesno = raw_input("Type 'YeSs' to submit the HiDRUS jobs.\t")

if(yesno!='YeSs'):
    quit()

#input for the command
cfile="./config.txt"
ens=range(1, 101)
pfile_prefix="./h2parms_erai_2008-14_"

\end{lstlisting}

\end{frame}
%
%
%
%
\begin{frame}[fragile]
\frametitle{Projection Run - PBS Job Script (Contd.)}
\begin{lstlisting}[language=Python]
# Loop over your jobs
for ens_id in ens:
    #make input parameters file name
    pfile=pfile_prefix+ str(ens_id)+ ".txt"
     
    # Customize your options here
    job_name = "h2era_"+str(ens_id)
    walltime = "10:00:00"
    mem="500MB"
    qtype = "normal"

    command = "./hidrus2 %s %d %s" %(cfile, ens_id, pfile)

\end{lstlisting}

\end{frame}
%
%
%
%
%
%
\begin{frame}[fragile]
\frametitle{Projection Run - PBS Job Script (Contd.)}
\begin{lstlisting}[language=Python]
 job_string = """#!/bin/bash
    #PBS -N %s
    #PBS -q %s
    #PBS -l walltime=%s
    #PBS -l mem=%s
    #PBS -P k10
    #PBS -l wd
    #PBS -m a
    #PBS -M bhupendra.raut@monash.edu
%s""" % (job_name, qtype, walltime, mem, command)
    
    # Open a pipe to the qsub command.
    output, input = popen2('qsub')
    # Send job_string to qsub
    input.write(job_string)
    input.close()
    
    #Print job and the system response as it's submitted
    print "\tsubmiting... "+job_name
    print output.read()
    time.sleep(0.5)
\end{lstlisting}

\end{frame}
%
%
%
%
%
%
\begin{frame}[fragile]
\frametitle{Projection Run - PBS Job Script (Contd.)}
\begin{lstlisting}[language=Python]
    #Print job and the system response as it's submitted
    print "\tsubmiting... "+job_name
    print output.read()
    time.sleep(0.5)
\end{lstlisting}

\end{frame}
%
%
%
%
%


\section{List of Programs and Scripts}
\begin{frame}{List of Programs and Scripts}
\begin{enumerate}
\item \textbf{mkCascParm\_dbz} : This C++ program computes STEPS parameters required to produce HiDRUS simulations.
\item \textbf{mk\_rainEvents\_lib.R} : This script creates library of rain-events and their parameters.

\item \textbf{add\_mslp\_uv\_rainlib.ncl} : This script is an example of adding relevant synoptic data to the rainlib for the better identification of the future days.
\item \textbf{mk-HiDRUS2-Parms.R} : This script taken in rain library files produced by above scripts and GCM/RCM data and assigns rainfall events to future days. Prints all the required parameters in an ascii file.
\item \textbf{hidrus2\_dbz} : This is the main simulation program which takes in the file generated by the \emph{mk-HiDRUS2-Parms.R} and produces downscaled rainfall realisation. You can run this program several time with the same input to get ensemble of realizations.
\item \textbf{nc2nc\_ents} : This program extracts the rainfall time-series at a given location from 100s of HiDRUS simulations and writes it in a netcdf file.

\end{enumerate}
\end{frame}

\begin{frame}{References}
\bibliographystyle{plain}
\bibliography{../references}
\end{frame}

\end{document}